\chapter{Conclusion and Future Work}
\label{ch:conclusion}

\section{Summary of Findings}
\label{sec:conc_summary}

\todo[inline,caption={Summarize the findings for RQ1, RQ2, and RQ3}]{Directly answer RQ1, RQ2, and RQ3 from Chapter~\ref{ch:introduction} using the verified results of Chapters~\ref{ch:mesa} and~\ref{ch:decisions}:
\begin{itemize}
    \item \textbf{RQ1:} Existing benchmark review frameworks agreed on 159 of 302 comparisons and penalized missing information 210 times versus 13 cited defects.
    \item \textbf{RQ2:} The 39-criterion \mesa{} rubric reached 88\% to 91\% verdict repeatability across Claude and Gemini, detected every planted defect, and produced 14 script-confirmed failed ratings on known defects.
    \item \textbf{RQ3:} On SymPy (0.70 requirement), 35 of 44 resolved on \swebench{} remained Inconclusive, strict review policies returned no decision due to alterable grading, independent patch checks kept 38 tasks (31 resolved, Inconclusive), and 44 post-cutoff SymPy issues gave 21 resolved (Reject).
\end{itemize}}

\section{Artifacts and Reproducibility}
\label{sec:conc_artifacts}

\todo[inline]{Summarize the public replication packages for both studies (mappings, rubric, frozen prompts/responses, decision engine, 88 frozen SymPy session records, unit tests, and verification scripts), linking to Appendix~\ref{app:decisions_artifacts}.}

\section{Future Work}
\label{sec:conc_future}

\todo[inline]{Outline next steps: (1) human studies evaluating \mesa{} usability, inter-rater agreement, and decision impact; (2) replicating the fail-closed requirement decision method across additional repositories, agents, and multiple runs per task; and (3) automating targeted benchmark-defect and patch-verification scripts.}
